\documentclass{article}
\usepackage{amsfonts,amsthm,amsmath}
\usepackage[mathscr]{euscript}
\usepackage{enumitem}

\setcounter{section}{0}
\renewcommand{\thesection}{\S\arabic{section}}
\renewcommand{\thesubsection}{\arabic{section}}

\newtheorem{thm}{Theorem}
\newenvironment{theorem}[1]{
	\renewcommand{\thethm}{#1}
	\begin{thm}
		}{\end{thm}}

\newtheorem{lma}{Lemma}
\newenvironment{lemma}[1]{
	\renewcommand{\thelma}{#1}
	\begin{lma}
		}{\end{lma}}

\theoremstyle{definition}
\newtheorem{ex}{Exercise}
\newenvironment{exercise}[1]{
	\renewcommand{\theex}{#1}
	\begin{ex}
		}{\end{ex}}

\title{Exercises 9.B}
\author{Connor Mooneyhan}
\date{March 22, 2024}

\begin{document}

\maketitle

\begin{ex}
	Suppose $S \in \mathcal{L}(\mathbb{R}^3)$ is an isometry. Prove that there exists a nonzero vector $x \in \mathbb{R}^3$ such that $S^2x = x$.
\end{ex}
\begin{proof}
	Since $\mathrm{dim}\,\mathbb{R}^3 = 3$ which is odd, given the characterization of an isometry from this section, there must be at least one $1\times1$ block consisting of the entry $1$ or $-1$.
	Let the basis element corresponding to that block be $e_j$.
	Then either \begin{equation*}
		S^2e_j = S(S(e_j)) = S(e_j) = e_j
	\end{equation*} or \begin{equation*}
		S^2e_j = S(S(e_j)) = S(-e_j) = e_j,
	\end{equation*} so $e_j$ satisfies the desired condition.
\end{proof}

\begin{ex}
	Prove that every isometry on an odd-dimensional real inner product space has $1$ or $-1$ as an eigenvalue.
\end{ex}
\begin{proof}
	Given the characterization of isometries from this section, the sum of the "dimensions" of the blocks must be equal to the dimension of the whole space.
	Thus, it cannot be made up of just $2\times2$ blocks, since then the sum would be even, so there must be at least one $1\times1$ block consisting of the entry $1$ or $-1$.
	Of course, that means that $1$ or $-1$ is an eigenvalue of the isometry as desired.
\end{proof}

\begin{ex}
	Suppose $V$ is a real inner product space. Show that \begin{equation*}
		\langle u + iv, x + iy \rangle = \langle u, x \rangle + \langle v, y \rangle + \left( \langle v, x \rangle - \langle u, y \rangle \right)i
	\end{equation*} for $u, v, x, y \in V$ defines a complex inner product on $V_C$.
\end{ex}
\begin{proof}
	We begin with definiteness and positivity.
	For $u + iv$ as given in the theorem statement, we have \begin{align*}
		\langle u + iv, u + iv \rangle & = \langle u, u \rangle + \langle v, v \rangle + \left( \langle v, u \rangle - \langle u, v \rangle \right)i \\
		                               & = ||u||^2 + ||v||^2 + (\langle v, u \rangle - \langle v, u \rangle )i                                       \\
		                               & = ||u||^2 + ||v||^2. \tag{1}
	\end{align*}
	\setcounter{equation}{1}
	This is clearly nonnegative for all values of $u$ and $v$, so positivity is satisfied.
	If either $u$ or $v$ is nonzero, its norm must also be nonzero, so (1) will also be nonzero.
	Conversely, if $u = v = 0$, then (1) is 0, so definiteness is satisfied.

	Moving to the "slot" properties, we must introduce yet another pair of vectors $p, q \in V$.
	The calculation for first-slot additivity is routine, if a bit horrific: \begin{align*}
		\langle u + iv, x + iy \rangle + \langle p + iq, x + iy \rangle =\  & \langle u, x \rangle + \langle p, x \rangle + \langle v, y \rangle + \langle q, y \rangle +  \\
		                                                                    & (\langle v, x \rangle + \langle q, x \rangle - \langle u, y \rangle - \langle p, y \rangle)i \\
		=\                                                                  & \langle u + p, x \rangle + \langle v + q, y \rangle +                                        \\
		                                                                    & \left( \langle v + q, x \rangle - \langle u + p, y \rangle \right)i                          \\
		=\                                                                  & \langle u + p + i(v + q), x + iy \rangle.
	\end{align*}
	First-slot homogeneity is no prettier.
	We define $a, b \in R$ and calculate: \begin{align*}
		(a + bi) \langle u + iv, x + iy \rangle =\  & (a + bi)\langle u, x \rangle + (a + bi)\langle v, y \rangle +                                       \\
		                                            & (a + bi)\left( \langle v, x \rangle - \langle u, y \rangle \right)i                                 \\
		=\                                          & \langle au, x \rangle + \langle bu, x \rangle i + \langle av, y \rangle + \langle bv, y \rangle i + \\
		                                            & (\langle av, x \rangle - \langle au, y \rangle)i - \langle bv, x \rangle + \langle bu, y \rangle    \\
		=\                                          & \langle au, x \rangle + \langle av, y \rangle + \langle bu, y \rangle - \langle bv, x \rangle +     \\
		                                            & (\langle bu, x \rangle + \langle bv, y \rangle + \langle av, x \rangle - \langle au, y \rangle)i    \\
		=\                                          & \langle au - bv, x \rangle + \langle av + bu, y \rangle +                                           \\
		                                            & \langle bu + av, x \rangle i + \langle bv - au, y \rangle i                                         \\
		=\                                          & \langle au - bv, x \rangle + \langle bu + av, y \rangle +                                           \\
		                                            & \langle bu + av, x \rangle i - \langle au - bv, y \rangle i                                         \\
		=\                                          & \langle au - bv + i(bu + av), x + iy \rangle                                                        \\
		=\                                          & \langle (a + bi)(u + iv), x + iy \rangle.
	\end{align*}

	Finally, for conjugate symmetry, we play with the inner product definition: \begin{align*}
		\langle u + iv, x + iy \rangle & = \langle u, x \rangle + \langle v, y \rangle + \left( \langle v, x \rangle - \langle u, y \rangle \right)i \\
		                               & = \langle x, u \rangle + \langle y, v \rangle + \left( \langle x, v \rangle - \langle y, u \rangle \right)i \\
		                               & = \langle x, u \rangle + \langle y, v \rangle - \left( \langle y, u \rangle - \langle x, v \rangle \right)i \\
		                               & = \overline{\langle x + iy, u + iv \rangle}.
	\end{align*}
\end{proof}

\begin{ex}
	Suppose $V$ is a real inner product space and $T \in \mathcal{L}(V)$ is self-adjoint.
	Show that $T_C$ is a self-adjoint operator on the inner product space $V_C$ defined by the previous exercise.
\end{ex}
\begin{proof}
	Let $u, v, x, y \in V$.
	Then \begin{align*}
		\langle T_C(u + iv), x + iy \rangle & = \langle Tu + iTv, x + iy \rangle                                                                  \\
		                                    & = \langle Tu, x \rangle + \langle Tv, y \rangle + \langle Tv, x \rangle i - \langle Tu, y \rangle i \\
		                                    & = \langle u, Tx \rangle + \langle v, Ty \rangle + \langle v, Tx \rangle i - \langle u, Ty \rangle i \\
		                                    & = \langle u + iv, Tx + iTy \rangle                                                                  \\
		                                    & = \langle u + iv, T_C(x + iy) \rangle
	\end{align*} as desired.
\end{proof}

\begin{ex}
	Use the previous exercise to give a proof of the Real Spectral Theorem via complexification and the Complex Spectral Theorem.
\end{ex}
\begin{proof}
	Given a self-adjoint $T \in \mathcal{L}(V)$ for some real inner product space $V$, the above exercise shows that $T_C \in \mathcal{L}(V_C)$ is self-adjoint, and in particular, it is normal.
	By the Complex Spectral Theorem, we can choose an orthonormal basis $e_1, \dots, e_n$ of $V_C$ consisting of eigenvectors of $T_C$ corresponding to eigenvalues $\lambda_1, \dots, \lambda_n$ (each not necessarily unique) respectively.
	Because $T_C$ is self-adjoint, the conjugate transpose of $\mathcal{M}(T_C)$ is equal to $\mathcal{M}(T_C)$.
	Since \begin{equation*}
		\mathcal{M}(T_C) = \begin{pmatrix}
			\lambda_1 &        &           \\
			          & \ddots &           \\
			          &        & \lambda_n \\
		\end{pmatrix},
	\end{equation*}
	\begin{equation*}
		\mathcal{M}(T_C*) = \begin{pmatrix}
			\overline{\lambda_1} &       &                      \\
			                     & \dots &                      \\
			                     &       & \overline{\lambda_n} \\
		\end{pmatrix},
	\end{equation*}
	and $\mathcal{M}(T_C) = \mathcal{M}(T_C*)$, we see that $\lambda_j = \overline{\lambda_j}$ for each $j \in \lbrace 1, \dots, n  \rbrace$.
	Thus each $\lambda_j$ is real.
	Write $e_j = u_j + iv_j$ for each $j$ where $u_j, v_j \in V$.
	Consider that each \begin{align*}
		Tu_j + iTv_j & = T_C e_j                          \\
		             & = \lambda_j e_j                    \\
		             & = \lambda_j (u_j + iv_j)           \\
		             & = \lambda_j u_j + i \lambda_j v_j.
	\end{align*}
	Thus $Tu_j = \lambda_j u_j$, so $u_j$ is an eigenvector of $T$ corresponding to $\lambda_j$.
	We proceed to show that the list $u_1, \dots, u_n$ is an orthonormal list.

	Given some $k \neq j$, consider that \begin{align*}
		0 & = \langle u_j + iv_j, u_k + iv_k \rangle                                                                       \\
		  & = \langle u_j, u_k \rangle + \langle v_j, v_k \rangle + (\langle v_j, u_k \rangle - \langle u_j, v_k \rangle)i
	\end{align*}
	So we get the two equations \begin{align*}
		0 & = \langle u_j, u_k \rangle + \langle v_j, v_k \rangle \\
	\end{align*} and \begin{equation*}
		\langle v_j, u_k \rangle = \langle u_j, v_k \rangle.
	\end{equation*}
	$\dots$
	\begin{align*}
		\langle u_j + v_j, u_k + v_k \rangle & = \langle u_j, u_k \rangle + \langle u_j, v_k \rangle + \langle v_j, u_k \rangle + \langle v_j, v_k \rangle \\
		                                     & = 2 \langle u_j, v_k \rangle                                                                                \\
	\end{align*}
  UNFINISHED
\end{proof}

\begin{ex}
	Give an example of an operator $T$ on an inner product space such that $T$ has an invariant subspace whose orthogonal complement is not invariant under $T$.
\end{ex}
\begin{proof}
	Define $T \in \mathcal{L}(\mathbb{R}^2)$ by the matrix (with respect to the standard basis) \begin{equation*}
		\begin{pmatrix}
			1 & 1 \\
			0 & 0
		\end{pmatrix}.
	\end{equation*}
	Then $T \lambda (1, 0) = \lambda T(1, 0) = \lambda (1, 0)$ so $\mathrm{span}\,(1, 0)$ is invariant.
  Meanwhile, the orthogonal complement of $\mathrm{span}\,(1, 0)$ is $\mathrm{span}\,(0, 1)$ but $T (0, 1) = (1, 0) \notin \mathrm{span}\,(0, 1)$ so $\mathrm{span}\,(0, 1)$ is not invariant under $T$.
\end{proof}

\end{document}
